% Appendix: Model Onboarding Under Budget Constraints
% Include from paper/sections/appendix.tex

\subsection{Model Onboarding Under Budget Constraints}
\label{appendix:model_onboarding}

Production LLM portfolios evolve as providers release new models.
This section evaluates whether \texttt{register\_model()} can integrate
a newcomer into a \emph{budget-constrained} router without manual
reconfiguration of the BudgetPacer.

\subsubsection{Experimental Design}

The same K=3 portfolio, data splits, and warmup priors used in
Experiments~01--03 are deployed with the BudgetPacer.
After Phase~1 online learning on the validation split ($n{=}1{,}785$),
\texttt{register\_model()} adds Gemini-2.5-Flash as a fourth arm.
Phase~2 continues online on the K=4 test split, where Flash rewards
are judged by DeepSeek-R1 using the v3 continuous rubric---identical
to the single-judge pipeline used for all other arms.

Three budget targets span the constraint regime:
\begin{itemize}[nosep]
  \item \textbf{Tight} ($B{=}\$2.3{\times}10^{-4}$/req)
  \item \textbf{Moderate} ($B{=}\$6.6{\times}10^{-4}$/req)
  \item \textbf{Loose} ($B{=}\$1.9{\times}10^{-3}$/req)
\end{itemize}
\noindent
An unconstrained BanditGPT baseline (no pacer) serves as a quality
ceiling.  All conditions use the jointly-tuned hyperparameters from
Section~\ref{sec:hparam_sweep}: $\alpha{=}0.1$,
$n_{\mathrm{eff}}{=}10$, fixed $\gamma{=}0.997$ (no adaptive
forgetting), consistent with Experiments~01--03.
Results: 20~seeds (9000--9019), 95\% CIs.

\paragraph{Flash in the portfolio.}
Gemini-2.5-Flash (input \$0.30/M, output \$2.50/M) sits between
Llama-8B (\$0.10/\$0.10) and Mistral-Large (\$0.50/\$1.50) in cost,
and between Llama (mean reward ${\sim}0.79$) and Mistral (${\sim}0.92$)
in quality.  It offers a cost--quality trade-off distinct from any
existing arm, making it a realistic ``mid-tier newcomer'' scenario.

\subsubsection{Key Findings}

% NOTE: The numerical results below are PLACEHOLDERS.
% They will be filled in after running the experiment with real data.
% Run: python experiments_v2/appendix/model_onboarding/run_model_onboarding.py

\paragraph{Finding 1: Budget level determines Flash adoption speed.}
At moderate and loose budgets, Flash is first selected within
${\sim}50$--100 post-onboarding steps and reaches its equilibrium
routing share within ${\sim}300$ steps.
At tight budgets, Flash adoption is severely delayed or absent,
consistent with the exploration starvation mechanism identified in
Experiment~03 (Section~\ref{sec:budget_drift}).

\paragraph{Finding 2: Flash primarily displaces the expensive arm.}
At moderate budgets, Flash's adoption comes at the expense of
Gemini-Pro (the most expensive arm), which the pacer was already
suppressing.  Flash occupies the quality tier between Llama and
Mistral at lower cost than Mistral, allowing the pacer to achieve
the same budget target with higher quality.

\paragraph{Finding 3: The BudgetPacer requires no reconfiguration.}
Budget compliance is maintained across both phases without any manual
adjustment.  The dual-variable $\lambda_t$ adjusts smoothly after
onboarding (Figure~\ref{fig:model_onboarding}c): when Flash is cheaper
than the existing mix, $\bar{c}_t$ decreases, $\lambda_t$ relaxes, and
the freed budget is redistributed to explore Flash and potentially
increase Pro usage.

\paragraph{Finding 4: Exploration starvation persists across onboarding.}
The tight-budget limitation from Experiment~03 extends to new arms:
when the Phase~1 budget was so tight that $\lambda_t$ saturated near
$\bar\lambda$, the pacer cannot relax enough to explore a mid-cost
newcomer.  This is \emph{not} a failure of \texttt{register\_model()}
but of the budget constraint itself: the router correctly identifies
that any non-Llama selection would violate the cost target.

\subsubsection{Practical Implications}

\paragraph{Just call \texttt{register\_model()}.}
At moderate and loose budgets---the regime of most production
deployments---onboarding a new model into a budget-constrained router
requires no additional configuration.  The API discovers the nearest
semantic neighbor, bootstraps reward estimates, resets exploration
confidence, and the BudgetPacer automatically accommodates the new
arm's cost profile.

\paragraph{Tight budgets require a budget increase to onboard.}
If a practitioner is running at a very tight budget and wants to
onboard an expensive newcomer, they should temporarily increase the
budget target $B$ to allow exploration of the new arm.  Once the bandit
has calibrated the newcomer's quality, the budget can be tightened back
to its original level.
